% Quantum-kernel components of QG-RSA (QKSR base + Q-BRIDGE incremental).
% Every quantity below corresponds to code in utils/quantum_kernel.py,
% utils/qbridge.py, utils/qbridge_calibration.py and models/RSIAT_adapter.py.

\section{Projected Quantum Kernels for Representation Steering}
\label{sec:quantum}

RSIAT~\cite{rsiat} learns a shared adapter on a frozen ViT-B/16 and shapes
its feature space with a cosine-similarity steering loss on the base task.
We ask whether a \emph{projected quantum kernel} (PQK)~\cite{huang2021power}
can replace this similarity with a non-linear, trainable metric, and
whether the same metric can tell, during incremental tasks, which old
classes a new sample is most likely to be confused with. The quantum
module is used only at training time; inference keeps the classical
ViT--adapter--cosine-classifier path unchanged.

\subsection{Projected Local Quantum Kernel}
\label{sec:pqk}

\subsubsection{Classical projection.}
A feature $\mathbf{z}\in\mathbb{R}^{768}$ is mapped to $q$ rotation angles
\begin{equation}
  \boldsymbol{\theta}(\mathbf{z}) = \pi \tanh\!\big(W\,\mathrm{LN}(\mathbf{z}) + \mathbf{b}\big)
  \in (-\pi,\pi)^{q},
  \label{eq:proj}
\end{equation}
where $\mathrm{LN}$ is a layer normalisation without affine parameters and
$W\in\mathbb{R}^{q\times 768}$. We use $q=8$ qubits.

\subsubsection{Circuit.}
The angles are loaded by single-qubit $R_y$ rotations and followed by $L=2$
variational layers, each consisting of trainable $R_y(\phi_{l,j})$ rotations
and a ring of CNOT gates $\mathrm{CNOT}_{j\to (j+1)\bmod q}$:
\begin{equation}
  |\psi(\mathbf{z})\rangle =
  \prod_{l=1}^{L}\Big[U_{\mathrm{ent}}\bigotimes_{j=1}^{q} R_y(\phi_{l,j})\Big]
  \bigotimes_{j=1}^{q} R_y\big(\theta_j(\mathbf{z})\big)\,|0\rangle^{\otimes q}.
  \label{eq:circuit}
\end{equation}
Since $R_y$ and CNOT are real, the $2^{q}=256$ amplitudes are real. The
circuit is simulated exactly and in batch as a differentiable statevector
in PyTorch, so no sampling noise is involved.

\subsubsection{Local reduced density matrices and kernel.}
Following~\cite{huang2021power}, we compare states through their one-qubit
reduced density matrices
$\rho_j(\mathbf{z}) = \mathrm{Tr}_{\bar{j}}\,|\psi(\mathbf{z})\rangle\langle\psi(\mathbf{z})|$
rather than through the global fidelity. The global fidelity is known to
concentrate as the number of qubits grows~\cite{thanasilp2024concentration}.
The distance and kernel are
\begin{equation}
  D(\mathbf{z},\mathbf{z}') = \sum_{j=1}^{q}
  \big\|\rho_j(\mathbf{z})-\rho_j(\mathbf{z}')\big\|_F^{2},
  \qquad
  k(\mathbf{z},\mathbf{z}') = \exp\!\big(-\gamma\,D(\mathbf{z},\mathbf{z}')\big).
  \label{eq:pqk}
\end{equation}
The bandwidth is bounded around a data-driven scale,
$\gamma = \gamma_0\,10^{\tanh(a)}$ with a trainable scalar $a$, where
$\gamma_0 = 1/\mathrm{median}\,D$ is computed once on 512 base-task training
features (off-diagonal pairs). Hence $\gamma$ stays within
$[\gamma_0/10,\,10\,\gamma_0]$, which prevents the kernel from collapsing to
the identity or to a constant matrix.

\subsection{QKSR: Quantum Kernel Steering on the Base Task}
\label{sec:qksr}

On the base task $t=0$, the PQK replaces the cosine similarity inside the
RSIAT steering loss. For a batch with labels $y$, let
$\mathcal{P}=\{(i,i'):y_i=y_{i'},\,i\neq i'\}$ and
$\mathcal{N}=\{(i,i'):y_i\neq y_{i'}\}$. Then
\begin{equation}
  \mathcal{L}^{Q}_{\mathrm{RS}} =
  \frac{1}{|\mathcal{P}|}\sum_{(i,i')\in\mathcal{P}}\big[1-k_{ii'}\big]_+
  + \alpha\,\frac{1}{|\mathcal{N}|}\sum_{(i,i')\in\mathcal{N}}\big[k_{ii'}-m_q\big]_+ ,
  \label{eq:qksr}
\end{equation}
where $k_{ii'}=k(\mathbf{z}_i,\mathbf{z}_{i'})$, $m_q=0.5$ and $\alpha=1$.
The base objective is
$\mathcal{L}_{\mathrm{cls}} + \lambda_{\mathrm{RS}}\min(1,e/E_w)\,\mathcal{L}^{Q}_{\mathrm{RS}}$
with a linear warm-up over $E_w$ epochs. The projector is trained with the
adapter learning rate, and the circuit parameters and $a$ with a ten times
smaller rate. After the base task, the PQK is frozen.

\subsection{Q-BRIDGE: Risk-Weighted Cross-Task Margins}
\label{sec:qbridge}

During incremental tasks $t>0$, RSIAT keeps the classification loss on the
new classes, the feature alignment
$\mathcal{L}_{\mathrm{align}}=\|\mathbf{z}-g(\mathbf{z}^{-})\|_2^2$ between
the current feature and the projected feature of the previous model
(mean-squared error, with $g$ the residual projector of RSIAT), and a
separation term between old prototypes and new samples. Q-BRIDGE replaces
the separation term with two margin losses and uses the frozen base PQK
only to \emph{weight} them.

\subsubsection{Cross-task margin.}
Let $s_{ic}$ be the cosine score of sample $i$ for class $c$, computed with
the old classifier rows detached so that the added losses cannot move old
decision boundaries directly. For each new-class sample $i$ and old class
$c$, the margin violation is
$v_{ic}=[\,m_x + s_{ic} - s_{iy_i}\,]_+$, and the loss is
\begin{equation}
  \mathcal{L}_{\mathrm{cross}} = \frac{1}{B\,C_{\mathrm{old}}}
  \sum_{i,c}\big(1 + g\,r_{ic}\big)\,v_{ic}^{2},
  \qquad
  r_{ic} = \sigma\!\big(\mathrm{softplus}(u)\,k(\mathbf{z}^{-}_i,\boldsymbol{\mu}_c)+b\big).
  \label{eq:cross}
\end{equation}
Here $\mathbf{z}^{-}_i$ is the previous model's feature, $\boldsymbol{\mu}_c$ is
the stored mean of old class $c$, and $g=0.5$. The risk $r_{ic}$ is
detached. Its monotone calibrator $(u,b)$ is fitted once by unweighted
binary cross-entropy on base-task training data only: prototypes, fitting
and validation pairs come from a stratified 60/20/20 split. The target is
whether the base classifier violates a margin of $0.1$ for that
(sample, class) pair. Because the calibrator is monotone, it preserves the
ranking of raw PQK similarities and only rescales them into a probability
of confusion. The calibrator is never refitted on incremental data.

\subsubsection{Old-anchor guard.}
To keep old classes from being absorbed by new ones, the projected old
prototypes $\mathbf{a}_c = g(\boldsymbol{\mu}_c)$ must score their own class
above every new class by a margin $m_{\mathrm{old}}$:
\begin{equation}
  \mathcal{L}_{\mathrm{guard}} = \frac{1}{C_{\mathrm{old}}}\sum_{c}
  \Big[\,m_{\mathrm{old}} + \max_{n\in\mathcal{C}_{\mathrm{new}}} s(\mathbf{a}_c,n)
  - s(\mathbf{a}_c,c)\,\Big]_+^{2}.
  \label{eq:guard}
\end{equation}
Both the anchors and the old classifier rows are detached, so this loss
only moves the new-class rows away from old prototypes.
The incremental objective is
$\mathcal{L}_{\mathrm{cls}} + \beta\,\mathcal{L}_{\mathrm{align}}
+ \lambda_x\mathcal{L}_{\mathrm{cross}} + \lambda_g\mathcal{L}_{\mathrm{guard}}$
with $m_x=m_{\mathrm{old}}=0.1$ and $\lambda_x=\lambda_g=0.1$. The drift
compensation of old statistics (SSCA) and the Gaussian classifier
alignment of RSIAT are kept unchanged.

\subsection{Controls for Attributing the Effect}
\label{sec:controls}

A gain from Eq.~\eqref{eq:cross} could come from the margins themselves,
from any similarity-based weighting, or specifically from the quantum
circuit. We therefore pre-register controls that differ from the full model
(R6, calibrated PQK) in one factor each:
\begin{itemize}
  \item \textbf{No separation:} CosFace over all seen classes with alignment only (R1).
  \item \textbf{Margins only:} uniform weights ($g=0$), without guard (R2) and with guard (R3).
  \item \textbf{Classical selectors:} cosine similarity (R4); an RBF kernel on the same
        learned projection $\boldsymbol{\theta}(\mathbf{z})$, without the circuit (R7);
        and a small MLP metric on $\boldsymbol{\theta}(\mathbf{z})$ (R8).
  \item \textbf{Uncalibrated PQK:} raw PQK risk (R5).
  \item \textbf{Information controls:} PQK risks shuffled among active pairs (R9) or replaced
        by their mean (R10). These keep the weight distribution but remove the
        pairing between risk and violation.
\end{itemize}
A quantum-specific contribution is claimed only if R6 outperforms R4, R7 and R8,
and R5 outperforms R9 and R10.

\subsubsection{Scope.}
The circuit has $q=8$ qubits and is simulated exactly, so it is classically
tractable by construction. We do not claim a quantum advantage. The
question studied here is whether a PQK, used as an inductive bias for a
learned metric, improves exemplar-free class-incremental learning relative
to matched classical metrics under the measurement protocol of
Section~\ref{sec:protocol}.
